% !TeX root = ../collection-solutions.tex

\titledquestion{Use it or lose it: taxing wealth instead of capital income}

\citet{guvenen2023use} study what happens when the United States replaces its capital income tax with a flat tax on wealth.

\begin{parts}
    \part[4] What is the key takeaway of the paper? Explain how the authors arrive at this conclusion: what kind of model do they use, how is it disciplined by data, and what is the main quantitative experiment behind the headline result?

    \begin{solutionorbox}[12cm]
        \emph{Takeaway} (2 points): When returns to wealth differ persistently across people, taxing wealth and taxing capital income are no longer equivalent. Replacing the US capital income tax with a revenue-neutral flat wealth tax raises output, wages and welfare (newborn welfare gain of $\approx 7\%$ in consumption-equivalent terms, with about two thirds of the population gaining). The optimal wealth tax is high ($\approx 3\%$), while the optimal capital \emph{income} tax would be a subsidy.

        \emph{How they get there} (2 points):
        \begin{itemize}
            \item A quantitative overlapping-generations model in which entrepreneurs differ in productivity (rate of return) and face collateral constraints, so wealth is not automatically allocated to its most productive use.
            \item The calibration matches US wealth concentration (top 1\% wealth share, Pareto tail, share of self-made billionaires); the model also reproduces the dispersion and persistence of returns measured in Norwegian administrative data \emph{without targeting them}.
            \item Main experiment: a revenue-neutral tax swap (capital income tax out, flat wealth tax of $\approx 1.2\%$ in). Capital rises by $\approx 16\%$, aggregate TFP by $\approx 2\%$, wages by $\approx 8\%$ — and most of the welfare gain comes through higher wages, which is why gains are broad-based even though wealth inequality \emph{rises}.
        \end{itemize}
        Grading: 1 point for the non-equivalence/efficiency-ranking insight, 1 point for the headline welfare result; 1 point for model + calibration, 1 point for the experiment with (approximate) magnitudes and the wage channel.
    \end{solutionorbox}

    \part[4] A common intuition is that a wealth tax punishes saving and should therefore shrink the economy. Explain intuitively the \emph{use-it-or-lose-it} mechanism by which the wealth tax raises productivity in this model.

    \begin{solutionorbox}[12cm]
        \begin{itemize}
            \item (1 point) A capital income tax is paid only on income \emph{generated}: productive entrepreneurs bear the entire burden, while idle wealth escapes taxation. A wealth tax is paid on the \emph{stock}: two people with the same wealth pay the same tax regardless of how productively they use it.
            \item (1 point) Switching to the wealth tax therefore broadens the tax base and shifts the burden from high-return to low-return wealth holders: ``use it or lose it''.
            \item (1 point) Because after-tax return \emph{differences} are preserved (rather than compressed by the income tax), behavior responds: high-return entrepreneurs save more, low-return holders run down their wealth. With collateral constraints, the wealth shed by unproductive holders is lent to productive but constrained entrepreneurs.
            \item (1 point) The same aggregate capital stock then produces more output: TFP rises, labor demand and wages rise, so the gains accrue broadly — even though wealth concentration at the top increases.
        \end{itemize}
    \end{solutionorbox}
\end{parts}
